% Generated draft: requires mathematical and source review.
\documentclass[11pt]{article}
\usepackage[margin=25mm]{geometry}
\usepackage{fontspec}
\setmainfont{Latin Modern Roman}
\usepackage{amsmath,amssymb,mathtools,unicode-math}
\setmathfont{Latin Modern Math}
\usepackage{enumitem,needspace,xurl,hyperref}
\setlength{\parindent}{0pt}
\setlength{\parskip}{5pt}
\setlength{\emergencystretch}{3em}
\newcommand{\N}{\mathbb{N}}
\newcommand{\Z}{\mathbb{Z}}
\newcommand{\Q}{\mathbb{Q}}
\newcommand{\R}{\mathbb{R}}
\newcommand{\C}{\mathbb{C}}
\newcommand{\F}{\mathbb{F}}
\newcommand{\A}{\mathbb{A}}
\newcommand{\PP}{\mathbb{P}}
\newcommand{\E}{\mathbb{E}}
\newcommand{\eps}{\varepsilon}
\newcommand{\dd}{\,\mathrm{d}}
\begin{document}
\section*{UnsolvedMath ID: 30005044}
\textbf{Source problem code:} catalogue entry
\subsection{Definitions and mathematical statement}
Let \(T\) be a Galton--Watson tree, and let \(\xi\) denote its offspring count: different vertices have independent offspring counts with the same distribution as \(\xi\). Assume that
\[
\mathbb E[\xi]=\infty.
\]
The supplied question is: \emph{Can the contact process with fitness on \(T\) exhibit explosion under this infinite-mean offspring assumption?}

This is an unresolved formulation rather than a fully specified probability statement. The supplied record does not define the fitness variables, their distribution or dependence, the infection and recovery rates, the initial infected configuration, or the intended meaning of explosion. In particular, it does not specify whether explosion concerns infinitely many transition events or infinitely many distinct infections in finite time. It also does not specify whether the question asks for an example of an offspring and fitness law, a conclusion for every infinite-mean offspring law, or a characterization of the laws permitting explosion. No such choices are imposed here. The probability measure, any conditioning on survival of the tree, and the intended probability quantifier also require clarification.

The original excerpt asks whether one can ``close the gap,'' but the supplied record does not identify the results delimiting that gap. Its supplied literature assessment treats explosion in this contact-process model as an open question and explicitly states that explosive branching-process examples do not resolve it. These are source qualifications, not independently verified mathematical conclusions.
\subsection{Short English statement}
The question asks whether a contact process with fitness on a Galton--Watson tree can explode when the mean offspring count is infinite. A precise formulation requires the fitness model, transition rates, initial state, definition of explosion, and intended quantifiers. The supplied open status requires independent review.
\subsection{Sources}
\begin{itemize}
\item Supplied local UnsolvedMath registry record, permanent ID \(30005044\), title \emph{Explosion with Infinite-Mean Offspring}, exact source problem code \texttt{catalogue entry}. Project registry references: \texttt{data/unsolvedmath/problems.json} and \texttt{catalog/problems.json}.
\item Supplied source citation: \emph{Population Dynamics and Statistical Physics in Synergy} (2023), Oberwolfach Reports, pp.~577--656, DOI \texttt{10.4171/owr/2022/12}. \url{https://doi.org/10.4171/owr/2022/12}. The record attributes the question to an Oberwolfach workshop report, gives proposed year 2022, and names no individual proposer. The excerpt does not identify a precise page or speaker.
\item Supplied literature reference: \emph{A special family of Galton Watson processes with explosions}. \url{https://arxiv.org/abs/1502.07538}. The supplied assessment cautions that explosive branching examples do not settle the contact-process question.
\item Literature checked at: \texttt{2026-08-22}. This exact date belongs to supplied dated curation, not an independently verified literature assessment. That curation labels the problem open and describes explosion with infinite mean offspring as an explicit open question in the source report.
\item The registry labels the statement \texttt{corrected\_verified} and reports that it was reconstructed and checked against the official Oberwolfach report. These are supplied verification claims; no independent source consultation or mathematical verification is asserted here.
\end{itemize}
\end{document}
